\section{Scale disjunctions and shared operating variables}\label{app:scaling}
Scaling gives an example where a representation can remove an artificial
gap. The arguments are affine interpolation and classical disjunctive
convexification \citep[Section 2, Theorem 2.1]{Balas1998}; we give them
explicitly to identify the role of shared intensive variables.
Let $E\subset\R^d\times\R^p$ be nonempty compact, with per-unit extensive
coordinates $\xi$ and intensive coordinates $T$. Define
\[
 \lambda\odot E=\{(\lambda\xi,T):(\xi,T)\in E\},\qquad
 F(\Lambda,E)=\{(\lambda,\lambda\xi,T):\lambda\in\Lambda,(\xi,T)\in E\}.
\]
The nonnegative scale set $\Lambda$ is nonempty compact, with extremes
$\lambda_-,\lambda_+$. At zero scale, $X=0$ and $T\in\operatorname{proj}_T E$;
a different off-state domain must be modeled explicitly.

\begin{proposition}[Endpoint hull and multiplicities]\label{prop:scale-hull}
With $C=\conv E$, the hull of $F(\Lambda,E)$ is the hull of just its two
extreme-scale slices, with $E$ replaced by $C$. It is exactly the projection
of
\begin{equation}\label{eq:scale-disaggregation}
 \begin{gathered}
 0\le\theta\le1,\quad\lambda=(1-\theta)\lambda_-+\theta\lambda_+,
 \quad X=X^-+X^+,\quad T=T^-+T^+,\\
 (X^-,T^-)\in(1-\theta)(\lambda_-\odot C),\qquad
 (X^+,T^+)\in\theta(\lambda_+\odot C),
 \end{gathered}
\end{equation}
where $0K=\{0\}$ for nonempty $K$. For $N\ge1$ identical units sharing
$T$, allow integer $0\le n\le N$, $X=\sum_{i=1}^n\xi_i$, and
$(\xi_i,T)\in E$. Their hull is the same as $\conv F(\{0,\ldots,N\},E)$.
Writing $P_T=\operatorname{proj}_T C$, it is
\begin{equation}\label{eq:multiplicity-hull}
 0\le n\le N,\qquad
 \exists T^+:\ (X,T^+)\in\frac nN(N\odot C),\quad
 T-T^+\in(1-n/N)P_T.
\end{equation}
\end{proposition}
\begin{proof}
For fixed $(\xi,T)$, the point $(\lambda,\lambda\xi,T)$ depends affinely
on $\lambda$. If the extremes differ, interpolate with
$\theta=(\lambda-\lambda_-)/(\lambda_+-\lambda_-)$; if they coincide,
there is only one slice. Every original point is therefore a combination
of endpoint points. A fixed-scale map is linear in $(\xi,T)$, so its
hull is that slice of $C$. Combining two compact convex slices gives
\eqref{eq:scale-disaggregation}, including zero weights.
For an integer $n>0$ and fixed $T$, a sum of $n$ points of $E_T$ lies in
$n\conv E_T$, and identical copies include $n E_T$. Equivalently, a
multiplicity point is the average of its $n$ same-scale points
$(n,n\xi_i,T)$, while every same-scale point is a feasible multiplicity
point. Their hulls therefore agree at fixed $n$ after convexification.
Taking the hull over $n$ and applying the endpoint result proves
\eqref{eq:multiplicity-hull}. At $n=0$ both models use the stated
projection off-state; $N=0$ consists only of that slice.
\end{proof}
If $C=\{(\xi,T):A\xi+BT\le b\}$ and $P_T=\{T:DT\le e\}$ are
polytopes, \eqref{eq:multiplicity-hull} becomes
\[
 AX+NB T^+\le bn,\qquad D(T-T^+)\le e(1-n/N),\qquad0\le n\le N.
\]
Compact conic-representable slices can also be homogenized at zero weight.
Indeed, if $K_0=\{z:\exists y,Az+By+b\in\mathcal K\}$ is nonempty compact,
then at weight $t>0$ division in $Az+Bw+tb\in\mathcal K$ gives $tK_0$.
At $t=0$ a feasible $(z,w)$ is a recession direction of the lifted set:
add any nonnegative multiple to a feasible lift. Its projected direction
must be zero by compactness of $K_0$. Conversely $z=w=0$ is feasible.
This checks zero weights even when the auxiliary lift is unbounded.
The aggregation mechanism also appears in \citet[Section 4.1, Lemma 3
and Theorem 3, under Assumption 1]{WuEtAl2026}.

For example, let per-unit load $\xi\in[\ell,u]$, $0\le\ell<u$, common
$T\in[a,b]$, $a<b$, and per-unit extensive output $\omega=\xi T$.
Then $C$ is the McCormick hull of this bilinear graph. Formula
\eqref{eq:multiplicity-hull} supplies its exact local aggregate hull with
one disaggregated intensive coordinate. Its endpoint vertices are the two
off-state points and the four corners at scale $N$. Besides
$n\ell\le X\le nu$, $0\le n\le N$, $a\le T\le b$ and the global
McCormick inequalities for $W=XT$ on $[0,Nu]\times[a,b]$, this hull has
\begin{equation}\label{eq:scale-bilinear-cuts}
 W\le bX+\ell[NT-Na-n(b-a)],\qquad
 W\ge aX+\ell[NT-Nb+n(b-a)].
\end{equation}
To check sufficiency without a facet calculation, introduce $S=NT^+$.
The scaled per-unit McCormick system is equivalent to
\[
 an\le S\le bn,\quad N T-b(N-n)\le S\le NT-a(N-n),
\]
\[
 \ell S+aX-a\ell n\le W\le\ell S+bX-b\ell n,
 \qquad uS+bX-bun\le W\le uS+aX-aun.
\]
Eliminating $S$ by comparing each lower and upper bound gives precisely
the displayed system (redundant comparisons follow from the stated
bounds). When $\ell=0$, use these inequalities directly without dividing
by $\ell$. Alternatively \eqref{eq:multiplicity-hull} itself is a complete
formulation at that boundary.

\subsection{When scale costs do and do not separate}
\begin{example}[A scale-cost gap on a convex per-unit set]
\label{ex:scale-cost}
Take $E=\{(\xi,T):\xi=T\in[0,1]\}$, $\Lambda=\{0,1,2\}$, and
$f(0)=f(2)=1$, $f(1)=0$. The point $(\lambda,X,T)=(1,1,1/2)$ is the
average of $(0,0,0)$ and $(2,2,1)$. Its true minimum convexified cost
is one although the separate lower scale-cost envelope has $\check f(1)=0$.
In any representation let $p_i$ be the mass at scale $i$. Mean scale one
gives $p_0=p_2$. On scales zero, one, and two, respectively, $X-T$ is
at most zero, zero, and one. Thus the target $X-T=1/2$ forces
$p_2\ge1/2$, and the cost is at least $p_0+p_2\ge1$. The endpoint
representation attains it. Scaling $f$ scales this error arbitrarily.
\end{example}
An operating cost must also be lifted before convexifying a nonconvex
per-unit set. With $E=\{-1,1\}$, scale one and cost $|\xi|$, its
convexified cost at $X=0$ is one, not $|0|$. Further,
$\lambda c(X/\lambda,T)$ need not be convex with an unscaled intensive
coordinate: the convex cost $c(\xi,T)=T^2$ gives $\lambda T^2$.

\begin{proposition}[Extensive-only perspective]\label{prop:scale-perspective}
Let $E\subset\R^d$ be nonempty compact, $c:E\to\R$ continuous, and
$\Lambda\subset[0,\infty)$ finite and nonempty. Let $h$ be the lower
boundary of $\conv\{(\xi,c(\xi)):\xi\in E\}$ and $\check f$ the lower
convex envelope of the finite scale-cost data. For $\lambda>0$, the exact
hull of $Z\ge\lambda c(\xi)+f(\lambda)$, $X=\lambda\xi$, is
\[
 \lambda\in\conv\Lambda,\quad X\in\lambda\conv E,\qquad
 Z\ge\lambda h(X/\lambda)+\check f(\lambda).
\]
If $0\in\Lambda$, its zero-scale slice is $X=0$, $Z\ge f(0)$.
\end{proposition}
\begin{proof}
For a representation with weights $\mu_j$ and positive mean scale
$\lambda$, the size-biased weights $\mu_j\lambda_j/\lambda$ sum to one.
They give per-unit cost at least $h(X/\lambda)$, while the original
weights give scale cost at least $\check f(\lambda)$. Conversely choose
attaining per-unit atoms with weights $\nu_i$ and attaining scale weights
$\mu_j$. Their product weights $\mu_j\nu_i$ attain both bounds and the
required $X$. Compactness gives finite attaining representations; adding
vertical slack gives every larger $Z$. At zero mean scale, nonnegativity
forces all positively weighted scales to be zero.
\end{proof}
This is the usual perspective argument after the cost has been included
in the per-unit graph. The multiplicity statement follows by treating
that cost as another additive extensive coordinate.

\begin{proposition}[Rectangularity criterion]\label{prop:scale-rectangularity}
Let $E\subset\R^d\times\R^p$ be nonempty compact and let the finite
catalogue $\Lambda\subset[0,M]$ satisfy $\{0,s,M\}\subseteq\Lambda$
and $0<s<M=\max\Lambda$. Put $C=\conv E$,
$Q=\operatorname{proj}_\xi C$, and $P_T=\operatorname{proj}_T C$.
Every scale-only cost separates, in the precise sense
\begin{equation}\label{eq:scale-cost-separation}
 \begin{split}
 &\conv\{(\lambda,X,T,Z):(\lambda,X,T)\in F(\Lambda,E),Z\ge f(\lambda)\}\\
 &\qquad=\{(\lambda,X,T,Z):(\lambda,X,T)\in\conv F(\Lambda,E),
                         Z\ge\check f(\lambda)\},
 \end{split}
\end{equation}
if and only if $C=Q\times P_T$. It suffices to test the single cost
$f(s)=0$, $f(\lambda)=1$ at every other catalogue scale.
\end{proposition}
\begin{proof}
The epigraph hull equals the compact hull of the cost graph plus the
nonnegative vertical ray. Hence all required minimum-cost representations
are attained; no closure operation is hidden. If the diagnostic cost
separates, every point of the scale-$s$ base-hull slice has a zero-cost
representation. All its constituents must have scale $s$, so this slice
is exactly $s\odot C$. Put $\alpha=s/M\in(0,1)$. For
$(\xi,T_1)\in C$ and $T_0\in P_T$, mixing the extreme slices gives
$(s,s\xi,\alpha T_1+(1-\alpha)T_0)$ in that slice. Thus the fibre of
$C$ at $\xi$ contains $\alpha T_1+(1-\alpha)T_0$. Iteration gives
$\alpha^jT_1+(1-\alpha^j)T_0$ in the fibre for every $j$. Closedness
puts $T_0$ in it, proving rectangularity.
Conversely, under rectangularity the base hull is
$\{(\lambda,X,T):0\le\lambda\le M,X\in\lambda Q,T\in P_T\}$.
For $\lambda>0$ choose $\xi=X/\lambda$ and use any attaining scale-cost
weights with the same $(\xi,T)\in C$ in every scale slice. Decompose
that point into $E$ within each slice; the resulting representation
attains $\check f(\lambda)$. At zero choose any $\xi\in Q$ and use
the zero slice. Every representation has cost at least $\check f$ by
convexity, proving equality. Universal separation implies the diagnostic
case, completing the equivalence.
\end{proof}
Without an interior scale, two endpoint weights are already fixed by
$\lambda$, so all scale-only costs separate even for nonrectangular $C$.
Affine scale costs always separate. The criterion concerns all scale-only
costs with the stated off-state domain, not arbitrary operating costs.

\subsection{Integral counts and aggregate rounding}
For nonempty compact extensive-only unit types $E_j\subset\R^d$ and an integral polytope
$P\subset\R_+^k$ of allowable counts, the exact aggregate hull is
\begin{equation}\label{eq:integral-count-hull}
 \{(n,X):n\in P,\ X\in\textstyle\sum_j n_j\conv E_j\}.
\end{equation}
Indeed it is convex and contains every integer multiplicity point.
For the reverse inclusion write $X=\sum_j n_j\xi_j$ with
$\xi_j\in\conv E_j$, decompose $n$ into integral vertices of $P$, and
keep these $\xi_j$ fixed. At each integral vertex the sum of scaled
convex hulls is the hull of the corresponding finite Minkowski sum.
This proves \eqref{eq:integral-count-hull}; no integer-decomposition
property beyond integrality of $P$ is needed.

For $n$ identical compact extensive sets in $\R^d$, the classical
Shapley--Folkman lemma \citep[Appendix 2, Lemma 2 and corollary]{Starr1969}
represents any $X\in n\conv E$ as a sum with at most $\min(d,n)$
constituents outside $E$. Replacing each exceptional constituent by any
point of $E$ changes the sum by at most $\min(d,n)\operatorname{diam}(E)$
in a chosen norm. Thus the same bound holds for the Hausdorff distance
between the two aggregate sets. An $L$-Lipschitz objective has gap at most
$L\min(d,n)\operatorname{diam}(E)$ when minimized over those sets alone.
A relative $O(d/n)$ bound additionally requires a positive normalization
of order $n$ and uniformly bounded $L$ and diameter.

All these identities concern the stated local or aggregate sets. Added
balance constraints can exclude the rounded point, and
$\conv(F\cap H)$ can be strictly smaller than $\conv F\cap H$.
Endpoint hull equality also does not license removing intermediate
catalogue choices from the feasible model. If the per-unit set depends
on scale, even the hull equality can fail: $\xi=\lambda$ at scales
$1,2,3$ gives $(\lambda,X)=(2,4)$ outside the segment joining $(1,1)$
and $(3,9)$.
